\section{Aspekty ekonomiczne}

% SLAJD 23
\begin{frame}{Nieruchomości – przepaść w cenie \cite{link4}}
    Wstęp: Mimo wyższych zarobków w mieście, realna jakość życia za te same pieniądze jest niższa.

    \vspace{0.5cm}

    \begin{columns}
        \begin{column}{0.48\textwidth}
            \textbf{\textcolor{red}{W mieście:}}
            \begin{itemize}
                \item Ogromna kwota za mały metraż.
                \item Przestrzeń ograniczona do mieszkania w bloku.
                \item Widok: często ściana sąsiedniego budynku.
            \end{itemize}
        \end{column}
        
        \begin{column}{0.48\textwidth}
            \textbf{\textcolor{green!60!black}{Na wsi:}}
            \begin{itemize}
                \item Cena mieszkania = Przestronny dom.
                \item Własny ogród i prywatność.
                \item Nieporównywalnie lepsze warunki dla rodziny.
            \end{itemize}
        \end{column}
    \end{columns}
\end{frame}

% SLAJD 24
\begin{frame}{Codzienne wydatki i styl życia \cite{link5}}
    \begin{columns}
        \begin{column}{0.48\textwidth}
            \textbf{\textcolor{red}{Koszty miejskie:}}
            \begin{itemize}
                \item Droższe usługi i parkingi.
                \item Wyższe ceny podstawowych zakupów.
                \item Pokusy drogich rozrywek.
                \item Środowisko sprzyjające konsumpcjonizmowi (obciążenie portfela).
            \end{itemize}
        \end{column}
        
        \begin{column}{0.48\textwidth}
            \textbf{\textcolor{green!60!black}{Oszczędności wiejskie:}}
            \begin{itemize}
                \item Dostęp do tańszej, lokalnej żywności.
                \item Niższe koszty usług.
                \item Darmowa rekreacja na łonie natury.
                \item Życie bez ciągłej presji finansowej.
            \end{itemize}
        \end{column}
    \end{columns}
\end{frame}

% SLAJD 25
\begin{frame}{Podsumowanie finansowe \cite{link5}}
    \begin{table}
        \centering
        \small
        \caption{Bilans zysków i strat}
        \begin{tabular}{|l|l|l|}
            \hline
            \textbf{Kategoria} & \textbf{Miasto} & \textbf{Wieś} \\ \hline
            Mieszkanie & Bardzo drogo, ciasno & Umiarkowanie, przestronnie \\ \hline
            Wydatki & Wysokie (pokusy) & Niskie (lokalne produkty) \\ \hline
            Jakość za cenę & Niska & Wysoka \\ \hline
        \end{tabular}
    \end{table}

    \vspace{0.5cm}

    \begin{block}{Wniosek: Siła nabywcza}
        Na wsi za każdą zarobioną złotówkę można kupić znacznie więcej przestrzeni, spokoju i lepszej jakości życia.
    \end{block}
\end{frame}

% SLAJD 26
\section{Podsumowanie}
\begin{frame}{Podsumowanie: Dwie filozofie życia}
    \begin{itemize}
        \item Wybór "Miasto vs Wieś" to starcie dwóch różnych filozofii życiowych.
        \item Miasto, mimo pozornej atrakcyjności, przegrywa w kluczowych kategoriach ludzkich potrzeb.
        \item Ucieczka z metropolii to nie kaprys, lecz \textbf{świadome poszukiwanie równowagi} w świecie zdominowanym przez pośpiech.
        \item Wieś daje szansę na odzyskanie kontroli nad czasem i budowanie autentycznych relacji.
    \end{itemize}

    \vspace{0.3cm}

    \begin{alertblock}{Nowa definicja sukcesu}
        Z dala od zgiełku sukcesu nie mierzy się pieniędzmi, lecz prawdziwą jakością życia, spokojem ducha i poczuciem przynależności.
    \end{alertblock}
\end{frame}